% SOURCE_AUTHORITY: sga5_fr_workpass.tex, lines 1372--1494 (LF-normalized unit; line 1440 comment hazard excluded from semantic text).
\section*{5. Dualidade local sobre as curvas}

\begin{statement}{Teorema 5.1}
Seja $X$ um esquema noetheriano regular de dimensão $1$. Suponhamos que $n$ seja primo com as características residuais de $X$. Então o teorema de pureza é válido sobre $X$ (3.1.4) e o complexo $A_X$ é dualizante (1.7).
\end{statement}

\begin{prooftext}[Demonstração]
Suponhamos demonstrada a pureza. Como $X$ satisfaz evidentemente a condição (reg) (3.1.1), e a condição $(\mathrm{cdloc})_n$ se verifica em virtude de (SGAA~X~2.2), conclui-se, por (3.2.1), que $\dim.\mathrm{q.inj}(A_X)=2$. Por outro lado, a condição (ii) de (1.7) se verifica por (3.3.1 b) (i)) (que foi demonstrada usando a pureza). Portanto, em virtude de (4.4), ficamos reduzidos a provar que $A_X$ é dualizante quando $X$ é estritamente local. Mas, para provar a pureza, que é uma questão local em uma vizinhança de um ponto fechado, também podemos supor que $X$ é estritamente local. Finalmente, ficamos reduzidos a provar o teorema quando $X$ é estritamente local.

Suponhamos, pois, que $X$ é estritamente local. Sejam $x$ seu ponto fechado, $(\pi=\car(k(x)))$, $\eta$ seu ponto genérico e $U=X-x=\Spec(k(\eta))$. Seja $G=\Gal(k(\ol{\eta})/k(\eta))$ o grupo fundamental de $U$. Sabe-se [CL, cap.~IV \S~2, exerc.~2] que se tem uma sequência exata canônica
\[
(5.1.1)\qquad 1\longrightarrow P\longrightarrow G\longrightarrow H\longrightarrow1,
\]
onde $H=\prod_{\ell\ne\pi}\Z_\ell(1)$, $\Z_\ell(1)=\varprojlim_{\nu} \mu_{\ell^\nu}$, e $P$ é um grupo pró-$\pi$. (Naturalmente, $H$ é isomorfo, mas não canonicamente, a $\prod_{\ell\ne\pi}\Z_\ell$.)

Sabe-se, por outro lado (SGAA~VIII~2), que a cohomologia étale de $U$ se interpreta como a cohomologia galoisiana de $G$, correspondendo os feixes (resp.~feixes construtíveis) de $A_U$-módulos aos $A$-módulos (resp.~$A$-módulos de tipo finito) sobre os quais $G$ atua continuamente.

Posto isso, os $\uH_x^i(A_X)$ são determinados pela sequência exata
\[
(5.1.2)\qquad 0\longrightarrow \uH_x^0(A_X)\longrightarrow \H^0(X,A_X)\longrightarrow \H^0(U,A_U)\longrightarrow \uH_x^1(A_X)\longrightarrow0,
\]
e os isomorfismos
\[
(5.1.3)\qquad \uH_x^i(A_X)\simeq \H^{i-1}(U,A_U)\qquad\text{para }i\geqslant2.
\]
Agora, trivialmente temos $\H^0(X,A_X)=\H^0(U,A_U)=A$, e por isso $\uH_x^0(A_X)=\uH_x^1(A_X)=0$.

Por outro lado, para um $G$-$A$-módulo (galoisiano) $M$, tem-se a sucessão espectral de Hochschild - Serre
\[
E_2^{pq}=\H^p(H,\H^q(P,M))\Longrightarrow \H^{*}(G,M).
\]
Como $n$ é primo com $\pi$ e $P$ é um pro-$\pi$-grupo, temos
\[
\H^q(P,M)=0\quad\text{para }q\ne0,
\]
e por isso $E_2^{pq}=0$ para $q\ne0$, de onde obtemos um isomorfismo canônico
\[
(5.1.4)\qquad \H^p(G,M)\simeq \H^p(H,M^P).
\]
Vê-se ainda, mediante um argumento análogo, que $\H^p(H,N)$ se identifica com $\H^p(\widehat{\Z}(1),N)$ quando se faz agir trivialmente sobre $N$ o fator $\widehat{\Z}(1)$. Tem-se, pois, finalmente, com a convenção anterior, um isomorfismo canônico
\[
(5.1.5)\qquad \H^p(G,M)\simeq \H^p(\widehat{\Z}(1),M^P).
\]
Ora, a cohomologia galoisiana de $\widehat{\Z}$ é bem conhecida (CG, cap.~I, p.~31): para um $G$-$A$-módulo $N$, tem-se
\[
(5.1.6)\qquad
\H^i(\widehat{\Z},N)=0\quad\text{para }i\geqslant2,
\qquad
\H^0(\widehat{\Z},N)=N^{\widehat{\Z}},
\qquad
\H^1(\widehat{\Z},N)=N_{\widehat{\Z}}.
\]
De (5.1.3), (5.1.5) e (5.1.6) deduz-se imediatamente que
\[
\H^i(U,A_U)=0\quad\text{para }i\geqslant2,
\]
e
\[
\H^1(U,A_U)\simeq A.
\]
Este último isomorfismo não é canônico, pois depende da escolha de uma identificação de $\widehat{\Z}(1)$ com $\widehat{\Z}$. No entanto, usando a sequência exata de Kummer, obtém-se (SGAA~XIX~1.3) um isomorfismo canônico (dado pela \textquotedblleft classe fundamental local\textquotedblright{})
\[
(5.1.7)\qquad \H^1(U,\mu_n)\simeq A.
\]
Fica demonstrada a pureza.

Demonstremos agora que $A_X$ é dualizante. Segundo (4.7.17), trata-se de demonstrar que, para todo $A_U$-módulo localmente constante construtível $M$, o emparelhamento
\[
(5.1.8)\qquad \H^i(U,M)\times\H^{1-i}(U,\uHom(M,A_U))\longrightarrow \H^1(U,A_U)\xrightarrow{\sim}(\mu_n)_x^{\otimes -1}
\]
é uma dualidade perfeita. Interpretaremos este emparelhamento em termos de cohomologia galoisiana. Utilizaremos o seguinte lema elementar.
\end{prooftext}

\begin{statement}{Lema 5.1.9}
Denotemos por $D=\Hom(\ ,A)$ o funtor dualizante de Pontryagin sobre a categoria dos $A$-módulos de tipo finito. Seja $M$ um $A$-módulo de tipo finito sobre o qual um grupo $G$ age $A$-linearmente.
\begin{enumerate}[label=(\roman*)]
\item Tem-se um isomorfismo canônico
\[
D(M^G)\simeq (DM)_G.
\]
\item Se $G$ é um grupo profinito que atua continuamente sobre $M$ e $G$ é de proordem prima com $n$, a seta canônica
\[
M^G\longrightarrow M_G
\]
é um isomorfismo.
\end{enumerate}
\end{statement}

\begin{prooftext}[Demonstração]
Por definição, temos
\[
M^G=\varprojlim_{s\in G}\Ker(M\xrightarrow{s-1}M),
\]
de onde
\[
D(M^G)\simeq \varinjlim_{s\in G}\Coker(DM\xrightarrow{s-1}DM)\simeq (DM)_G,
\]
que demonstra (i).

Demonstremos (ii). Seja $M'$ o $A$-módulo definido pela sequência exata
\[
0\longrightarrow M'\longrightarrow M\longrightarrow M_G\longrightarrow0.
\]
Tem-se $\H^1(G,M')=0$ porque $G$ é de proordem prima com $n$. A sequência exata de cohomologia implica então que a seta $M^G\to M_G$ é um epimorfismo. Mas, aplicando o funtor $D$ e levando em conta (i), deduz-se que é um monomorfismo, o que termina a demonstração.
\end{prooftext}

Aplicando (5.1.9) a $M$ e $P$, obtém-se que $\uHom(M,A)^P$ se identifica canonicamente com $\uHom(M^P,A)$. Se se escolher um isomorfismo de $\widehat{\Z}(1)$ sobre $\widehat{\Z}$, o emparelhamento (5.1.8) escreve-se, levando em conta (5.1.5),
\[
(5.1.8)'\qquad \H^i(\widehat{\Z},M^P)\times\H^{1-i}(\widehat{\Z},D(M^P))\longrightarrow\H^1(\widehat{\Z},A)\simeq A.
\]
De (5.1.6) e (5.1.9(i)) resulta imediatamente que (5.1.8)' é um emparelhamento perfeito. Isto termina a demonstração de (5.1).

\section*{Bibliografia}

\begin{description}[leftmargin=2.5cm,style=nextline]
\item[{[2]}] Hironaka, H. Resolution of singularities of an algebraic variety over a field of characteristic zero, Ann. of Math., vol. 79 (1964), p.~109.
\item[{[1]}] Abhyankar, S. Local uniformization on algebraic surfaces over ground fields of characteristic $p\ne0$, Ann. of Math., vol. 63 (1956), pp.~491-526, e Resolution of singularities of arithmetical surfaces, Arithmetical Algebraic Geometry, Harper \& Row, Nova York (1965).
\item[{[H]}] Hartshorne, R. Residues and duality, Lecture Notes n.\textsuperscript{o}~20, Springer (1966).
\item[{[SGAA]}] Séminaire de Géométrie Algébrique de l'IHES, Cohomologie étale des schémas, por M.~Artin e A.~Grothendieck, 1963-64.
\item[{[CL]}] Serre, J.-P. Corps locaux, Hermann (1962).
\item[{[CG]}] Serre, J.-P. Cohomologie galoisienne, Lecture Notes n.\textsuperscript{o}~5, Springer (1964).
\end{description}
